% Privacy-first section 6: what decides whether a memorised record can be reached.
% Budget 1.5 pages including Table~\ref{tab:slots} and Figure~\ref{fig:slots}.
% Numbers: outputs/kcurve300 (window 300, k=1..6) replicated by outputs/kcurve (window 288);
% DNABERT family from outputs/dnabert_long/*_human_s0, re-measured 23 Oct. Notes 3.58(f), 3.59.
\section{What decides whether a leak can be reached}
\label{sec:slots}

Released models differ enormously in how much comes back out of them, and the obvious explanation
--- that some memorise more than others --- is wrong. This section shows that what varies is not how
much is stored but how many pieces have to be correct simultaneously to get it out, and that this is
set by the tokenizer.

\paragraph{The same knowledge, 92\% or 0\% out.}
Hold the backbone, the corpus, the optimiser and the schedule fixed, and vary only how the DNA is
split. Hiding 48 nucleotides then occupies $48/k$ token slots, and every slot has to be right at the
same time for the record to come back whole. Table~\ref{tab:slots} reads each run at matched held-out
loss rather than at a matched epoch, because the tokenizers converge at very different rates: at
round~30 every $k\ge3$ has saturated at 100\% and the comparison says nothing at all.

Two things in that table carry the claim. The extraction rate is monotone in the slot count with no
exception across six cut sizes, and the loss it costs to reach full extraction falls monotonically in
the same direction: 4.07\,\bits{} at 16 slots, 2.99 at 12, 2.63 at 9.6, 2.47 at 8, and never reached
at 24 or 48. A coarse tokenizer does not merely leak more. It leaks more cheaply, in the same
currency \S\ref{sec:leak} prices released checkpoints in.

\begin{table}[t]
\centering
\caption{From-scratch backbone, next-token objective, three seeds, uniform DNA, read at matched
held-out loss. \emph{Slots} is $48/k$, the number of token positions the hidden 48 nucleotides
occupy. \emph{Cost of 100\%} is the held-out loss at which extraction first reaches saturation.
Never-trained controls read $0.00$ in every tokenizer at every snapshot.}
\label{tab:slots}
\begin{tabular}{lcccccc}
\toprule
tokenizer & slots & @2.05 & @2.10 & @2.20 & @2.40 & cost of 100\% \\
\midrule
character        & 48  & ---  & ---  & ---  & ---  & never \\
\kmer{2}         & 24  & 0\%  & 0\%  & 0\%  & 0\%  & never \\
\kmer{3}         & 16  & 0\%  & 0\%  & 0\%  & 2\%  & 4.07 \\
\kmer{4}         & 12  & 0\%  & 0\%  & 2\%  & 12\% & 2.99 \\
\kmer{5}         & 9.6 & 0\%  & 2\%  & 24\% & 61\% & 2.63 \\
\kmer{6}         & 8   & 1\%  & 17\% & 67\% & 92\% & \textbf{2.47} \\
\bottomrule
\end{tabular}
\end{table}

The character rows are empty because the character model never reaches the comparison points: its
held-out loss stays between 2.00 and 2.04 for thirty epochs and its extraction rate stays at zero
throughout. It does not resist extraction after overfitting; it does not overfit far enough to be
asked.

\paragraph{The effect is not an artefact of the window.}
A reader may reasonably suspect a geometry that was chosen to make $48/k$ come out neatly. We
repeated the whole sweep at a second window length, 300\nt{} instead of 288, which changes the number
of host positions, the number of training tokens per window and the offsets every canary lands on.
Where the two windows share a tokenizer they agree closely: at matched loss $2.40$, \kmer{3} reads
2\% against 3\%, \kmer{4} 12\% against 11\%, \kmer{6} 92\% against 91\%, and the cost of full
extraction is 4.07/2.99/2.47 against 4.15/3.01/2.47. Two independent sets of eighteen and fifteen
runs, at different geometries, trace the same curve.

\paragraph{A prediction that could have falsified this, and did not.}
The slot account makes a claim about released models that is easy to check and awkward to explain
away if it fails. DNABERT ships four checkpoints differing only in $k$ \citep{ji2021dnabert}, which
looks like the ideal natural experiment --- but its $k$-mers \emph{overlap}, advancing one nucleotide
at a time, so the slot count stays near 48 whatever $k$ is. The slot account therefore predicts no
trend in that family at all, in contrast to the strong trend it predicts for non-overlapping
tokenizers. We recorded that prediction before measuring.

Fine-tuned on chromosome~22 with real haplotype canaries and read with a nucleotide-level decoder
matched to overlapping spans, the four checkpoints recover 2\%, 0\%, 4\% and 0\% of records planted
sixteen times, for $k=3,4,5,6$ respectively. There is no trend, as predicted.

\paragraph{The flat row is not ignorance.}
It would be a misreading to conclude that the overlapping family has memorised nothing. Asked for any
single nucleotide of a planted record with all the others supplied, those same four checkpoints
answer correctly 0.99--1.00 of the time, and the implied ceiling on whole-record recovery runs from
0.70 to 0.87. The records are in there, close to completely; 48 simultaneous slots is simply a wall
that our search cannot climb. This is the mechanism of \S\ref{sec:slots} appearing in a second place:
the same quantity that makes a 6-mer model leak at 2.47\,\bits{} makes an overlapping model look
clean while holding the same content.

Stated carefully, then, the finding is not that finer tokenisation memorises less. Both ends of
Table~\ref{tab:slots} learn a planted record to the same per-nucleotide accuracy. Coarser
tokenisation makes what was learned \emph{reachable}, and that is a different sentence to hand a
practitioner, with a different remedy attached.
